\subsection{Rational normalization: the concurrent cubic and the cusp branch}

\begin{statement}{C.2.RAT.1}{Proved}{Euler signs for two of the rational boundary classes}
Let $f:Y\to B$ be a proper holomorphic map from a smooth compact
complex threefold with $b_2(Y)=0$, with connected fibres and smooth
general fibre a complex two-torus. Let $mS$ be a full fibre with
irreducible nonnormal reduced support. Write
$h=\eta\pi:Z\to T\to S$ for the minimal resolution of its
normalization and suppose that $Z$ is rational. Use the effective
divisor $K_Z+D=0$ of A.5.16.

If a contraction to $\mathbf P^2$ sends $D$ to three concurrent
reduced lines, then $e(S)\ge1$. More precisely,
$e(S)=-1+N_0+e(\Gamma_{\rm red})\ge1+N_0$, with $N_0$ as in
(A.RT.12).

Consider also a specified morphism $\rho:Z\to\mathbb F_n$,
$n\ge1$, which is a sequence of point blowups in reverse and has
\[
 \rho_*D=2C_0+\sum_{j=1}^r a_jF_j,\qquad
 C_0^2=-n,\qquad
 a_j\ge1,\quad \sum_j a_j=n+2,                  \tag{C.RA.1}
\]
where the $F_j$ are distinct ruling fibres. Let $C$ be the
strict transform of $C_0$. Suppose $h(C)$ is a curve and the
generic transverse singularity there is a cusp $A_2$.
Then $r\le3$ and $e(S)\ge1$. If $h(C)$ is smooth, then
$r\le2$. These assertions allow singular normalization,
arbitrary vertical conductor multiplicities and a nonreduced
conductor-image scheme.
\end{statement}
\begin{proof}
Put $G=\Gamma_{\rm red}$. A.5.16 has already proved
\[
 e(S)=e(Z\setminus D_{\rm red})+e(G).             \tag{C.RA.2}
\]
In either boundary class under consideration, $D_{\rm red}$
is simply connected, all its components are rational, and
(A.RT.18)--(A.RT.20) give
\[
 \pi_1(S)=\pi_1(G)\in\{1,\mathbb Z\}.
 \tag{C.RA.3}
\]
Write $n_G$ for the number of irreducible components of $G$
and $\beta=b_1(G)\le1$. The normalization maps from their
rational normalizations are homeomorphisms except for the
identification of finitely many points. Consequently
\[
                   e(G)=n_G+1-\beta\ge n_G\ge1.
 \tag{C.RA.4}
\]
The exact two-branch incidence argument (A.RT.21) strengthens
this to $e(G)\ge2$ in the present boundary classes.
For three concurrent lines, (A.RT.12)--(A.RT.13) therefore
prove the first assertion.

We prove the cusp assertion. Put $H=h(C)_{\rm red}$.
The transverse $A_2$ curve has a single normalized branch.
Thus $C\to H$ is the normalization map of degree one, and
no other strict conductor component maps onto $H$.
In particular $H^\nu=C=\mathbf P^1$. The restriction
$H^1(G,\OO_G)\to H^1(H,\OO_H)$ is surjective, since its
kernel has one-dimensional support and no second cohomology.
A.5.16 therefore gives
\[
              \delta:=p_a(H)\le2.              \tag{C.RA.5}
\]
Since $H^\nu$ is rational, $\delta$ is the sum of its local
delta invariants. Its topological incidence graph is a
subgraph of the incidence graph of $G$, so $b_1(H)\le1$.
If $s$ denotes the number of points of $C$ above singular
points of $H$, normalization of the reduced curve gives
\[
 s=\sum_{p\in\operatorname{Sing}H}r_p
   =b_1(H)+\#\operatorname{Sing}H\le1+\delta
 \quad(\delta>0),\qquad s=0\quad(\delta=0).       \tag{C.RA.6}
\]

We need the precise low-delta curve analysis, including its
two possible non-lci cases. A reduced curve germ of delta
invariant at most two has at most three branches.
Three branches require delta two, and the ring is exactly
the subring of its normalization consisting of triples
with equal constant value. This follows by equality of
colengths, both equal to two. Its topological branch
identification alone contributes two loops to a globally
irreducible curve, so $b_1(H)\le1$ excludes it here.

For two branches, the exact sequence for their union gives
$\delta=\delta_1+\delta_2+\ell(\OO_{H_1\cap H_2})$.
Two smooth branches therefore give a node, or, at delta
two, the fibre product matching their constant and linear
terms. The latter has the plane tacnode ring
$\mathbb C\{x,y\}/(y^2-x^4)$.
The only remaining two-branch possibility is a cusp and
a transverse smooth branch, whose ring is
$\mathbb C\{t^2,t^3\}\times_{\mathbb C}\mathbb C\{u\}$.
In embedding dimension three its ideal can be written
\[
                  I=(xz,yz,y^2-x^3).            \tag{C.RA.7}
\]
For one branch, the delta invariant equals the number of
gaps of its valuation semigroup. With at most two gaps,
the singular possibilities are
$\langle2,3\rangle$, $\langle2,5\rangle$, and
$\langle3,4,5\rangle$. The first two are plane cusps.
For the last, the conductor is $t^3\mathbb C\{t\}$,
and colength two makes the ring
$\mathbb C+t^3\mathbb C\{t\}$ exactly. Its ideal is
\[
 I=(g_1,g_2,g_3)
   =(y^2-xz,\ yz-x^3,\ z^2-x^2y),\qquad
 (x,y,z)=(t^3,t^4,t^5).                         \tag{C.RA.8}
\]
For completeness, the asserted plane form at semigroup
$\langle2,5\rangle$ follows by making an order-two
generator equal to $t^2$ by a change of parameter.
Its conductor is then $t^4\mathbb C\{t\}$, and the ring
is $\mathbb C+\mathbb Ct^2+t^4\mathbb C\{t\}$.
These conductor descriptions also follow in the completed
local rings by eliminating successive higher terms; finite
normalization and faithful flatness of completion give the
same conductor in the convergent rings. Thus there are no
unlisted cases. A germ of embedding dimension two is lci
in the smooth threefold: choose one regular transverse
ambient coordinate in its ideal and then its single plane
equation. We now exclude (C.RA.7)--(C.RA.8) for our cusp
image.

For (C.RA.8), put
\[
 F=x^5+xy^3+z^3-3x^2yz .
\]
The symbolic square is $I^{(2)}=I^2+(F)$, with
\[
 xF=g_2^2-g_1g_3,\quad
 yF=g_2g_3+xg_1^2,\quad
 zF=g_3^2+xg_1g_2.                              \tag{C.RA.9}
\]
The saturation computation $I^2:x^\infty=I^2+(F)$,
including the complete exact Groebner certificate, is
given in \path{code/A_low_delta_cusp.py} and its output.
It computes the symbolic square because the curve is lci
away from the origin and $x$ vanishes on it only at the
origin. The same equality in the analytic local ring
follows by flatness of the convergent extension and the
displayed finite saturation.

Let $N$ be the dual of the torsion-free pullback to $C$ of
the ambient conormal module of $H$. At (C.RA.8) the
normalized conormal basis and its dual lifts are
\[
\begin{gathered}
 a=dg_1,\qquad b=dg_2+t\,dg_1,\qquad
 dg_3=-2t^2a+tb,\\
 A=-t^{-3}\partial_z,\qquad
 B=(t^{-5}\partial_y+2t^{-4}\partial_z)/3 .
\end{gathered}                                  \tag{C.RA.10}
\]
Their wedge is $a\wedge b=t^6\,\iota_{\gamma'}(dx\wedge
dy\wedge dz)$, where $\gamma=(t^3,t^4,t^5)$.
Thus the determinant formula is
$\det N=K_C\otimes h_C^*K_Y^{-1}(6p)$.
Since every ambient line has degree zero on every compact
curve when $b_2(Y)=0$, its degree here is four.

Write the local equation of $S$ as $cF+f_2$, $f_2\in I^2$.
Its normalized quadratic form is the sum of a regular
quadratic form and $c$ times
\[
 q_F=\frac3t a^2-\frac3{t^2}ab+\frac1{t^3}b^2,
 \qquad \det(q_F)=\frac3{4t^4}.                  \tag{C.RA.11}
\]
If $c(0)\ne0$, the determinant has an uncancellable
pole of order four: a cross term with the regular form
has pole at most three. The generic cusp requires rank
one, so $c(0)=0$. Equations (C.RA.9) then give
$cF\in I^2$, and the quadratic form is regular.

The intrinsic cubic on its radical has pole order at most
ten at $p$. Here is a direct bound for arbitrary coefficients,
not only for the polynomial $F$. In the lifts (C.RA.10),
the differentials $dg_i$ are regular, the second derivatives
$D^2g_i(Au+Bv)^2$ have poles at most ten, and a derivative
of a holomorphic coefficient has pole at most five.
The formula for $D^3(b_{ij}g_i g_j)$ therefore has pole
at most ten. Terms containing an undifferentiated $g_i$
vanish on the curve. The exact derivatives, with equality
in the pole-ten bound, are in the cited computation.
The radical cubic is independent of a tangent component
of the lift: differentiating the identically zero quadratic
form on its radical along the smooth part of the curve
gives this invariance, and equality extends meromorphically.

The same argument is shorter for (C.RA.7). Put
$g_1=xz$, $g_2=yz$, $g_3=y^2-x^3$. Then
$I^{(2)}=I^2+(z g_3)$ and
\[
 x(zg_3)=g_1g_3,\quad y(zg_3)=g_2g_3,\quad
 z(zg_3)=g_2^2-xg_1^2.
 \tag{C.RA.12}
\]
On the cusp branch $(t^2,t^3,0)$ use $a=dg_1$, $b=dg_3$.
The quadratic form of $z g_3$ is $t^{-2}ab$, whose
determinant has pole four and cannot be cancelled by a
regular quadratic form. Thus its coefficient again
vanishes at the origin and the surface equation belongs
to $I^2$. The determinant corrections are four on the
cusp branch and two on the smooth branch $(0,0,t)$.
The total determinant degree is again four.
Dual lifts on the cusp branch are
$t^{-2}\partial_z$ and $\frac12t^{-3}\partial_y$;
on the smooth branch they are
$t^{-1}\partial_x,t^{-1}\partial_y$.
The same product derivative formula bounds the cubic
poles by six and two respectively, of total degree eight.

In either non-lci case let $R\subset N$ be the saturated
radical of the regular rank-one quadratic form and let
$Q=N/R$. Its nonzero square coefficient is a section of
$Q^{-2}\otimes M$, where $M=\OO_Y(S)|_C$ has degree
zero. Hence $\deg Q\le0$ and $\deg R=4-\deg Q\ge4$.
The nonzero intrinsic cubic of a generic $A_2$ is a
section of $R^{-3}\otimes M$ with poles of total degree
at most ten, or eight. It would give respectively
$3\deg R\le10$ or $3\deg R\le8$, both impossible.
This excludes the two non-lci cases, without assuming
that the conductor-image scheme is reduced.

We may consequently use the lci conormal bundle along $H$.
Let $\Xi$ be its normalization conductor divisor, of degree
$2\delta$. The determinant formula and the
conductor-cleared intrinsic cubic from A.5.5 give
\[
\begin{gathered}
 \det N=K_C(\Xi)\otimes h_C^*K_Y^{-1},\\
 q\in H^0(C,Q^{-2}\otimes M),\qquad
 k\in H^0(C,R^{-3}\otimes M(2\Xi)),\\
 0\ne q^3k^2\in
 H^0\bigl(C,(\det N)^{-6}\otimes M^5(4\Xi)\bigr),\\
 \deg\bigl((\det N)^{-6}\otimes M^5(4\Xi)\bigr)
                         =12-4\delta .
\end{gathered}                                  \tag{C.RA.13}
\]
The section is nonzero because the generic transverse
singularity is a cusp.

We give the local comparison that strengthens the zero
count at a smooth point of $H$. A simple zero of $q$ at
which $k$ is nonzero is impossible in a torus smoothing.
Choose coordinates with $H=\{x=y=0\}$, radical direction
$\partial_x$, and coordinate $z$ along $H$.
The nonzero cubic permits Weierstrass preparation in $x$:
\[
 f_S=U\bigl(x^3+a_2(y,z)x^2+a_1(y,z)x+a_0(y,z)\bigr).
 \tag{C.RA.14}
\]
The critical-curve equations and the rank-one Hessian
along it give $a_2\in(y)$ and $a_1,a_0\in(y^2)$.
After depressing the cube by $x'=x+a_2/3$, this is
\[
 f_S=U'\bigl(x'^3+y^2(B_1(y,z)x'+B_0(y,z))\bigr).
 \tag{C.RA.15}
\]
The simple zero of the Hessian coefficient says
$\partial_zB_0(0,0)\ne0$. Thus
$z'=B_1(y,z)x'+B_0(y,z)$ is an actual holomorphic
coordinate with $x',y$, and the reduced equation is a
unit times $x'^3+y^2z'$. The full fibre equation is a
unit times its $m$th power. A.5.15 excludes precisely
this Hirzebruch germ, including that unit and multiplicity.

Every point of $C$ above the smooth part of $H$ where a
vertical conductor component meets it has $q=0$.
Indeed a nonzero Hessian would, by A.5.17, give a constant
transverse cusp product with smooth normalization and pure
conductor $2C$; its minimal normalization resolution could
have no such vertical component. At a zero of $q$ the
cubic is holomorphic. If it also vanishes, the section
$q^3k^2$ has zero order at least five. If it does not
vanish, (C.RA.14)--(C.RA.15) force the zero of $q$ to
have order at least two, and the order is at least six.

Each of the $r$ original ruling fibres in (C.RA.1) has
a positive-boundary total-transform component meeting $C$,
at a different point of $C$. This remains true at each
blowup: a blowup at its intersection with the coefficient-two
section inserts an exceptional curve of positive coefficient.
Using (C.RA.6) and the degree in (C.RA.13) now gives
\[
 \begin{array}{c|c|c|c}
 \delta& s\text{ at most}&
 \text{bad points above }H_{\rm reg}\text{ at most}&r\text{ at most}\\ \hline
 0&0&2&2\\
 1&2&1&3\\
 2&3&0&3 .
 \end{array}                                    \tag{C.RA.16}
\]
This proves the asserted bound on $r$.

Finally (A.RT.12) and (A.RT.16) give
$e(Z\setminus D_{\rm red})=2-r+N_0$.
The forest argument of (A.RT.21) applies to this entire
star class, whether its minimal complement is negative
or not, and gives $e(G)\ge2$.
Together with $r\le3$, (C.RA.2) now gives directly
\[
                 e(S)\ge4-r+N_0\ge1+N_0.
 \tag{C.RA.17}
\]
\end{proof}

The two classes in C.2.RAT.1 are proved without a restriction
on nilpotent conductor structure. The other horizontal
coefficient-two types and contraction of that horizontal
section are not included in this statement.

\begin{statement}{C.2.RAT.2}{Proved}{A rational double cover of a smooth tacnode image}
Retain the torus-fibre, rational-normalization and specified
ruled-boundary morphism of C.2.RAT.1. Suppose that $C$ is not
contracted, that $H=h(C)_{\rm red}$ is a smooth rational curve,
that the generic transversal along $H$ is an ordinary tacnode,
and that $C\to H$ has degree two. Then $r\le2$ and
\[
                         e(S)\ge2+N_0.
\tag{C.RA.18}
\]
The assertion allows arbitrary vertical conductor and singular
normalization away from the generic points of $C$.
\end{statement}
\begin{proof}
The two branches of a generic tacnode are both supplied by
$C\to H$; no other strict conductor component can dominate $H$.
Apply the exact normal-jet construction (A.T.5)--(A.T.7) on
$H\simeq\PP^1$. Here the curve conductor divisor $\Xi$ is zero.
The quotient $Q$ of its rank-one Hessian has degree $q_0\le0$.
The quadratic coefficient $a$ and the quartic invariant
$\Delta=ae-b^2/4$ are nonzero sections with
\[
 A:=\deg\operatorname{div}(a)=-2q_0,\qquad
 \deg\operatorname{div}(\Delta)=8+2q_0=8-A.
\tag{C.RA.19}
\]
In particular $A$ is a nonnegative even integer. These are
sections of the actual lines $M\otimes Q^{-2}$ and
$(M\otimes Q^{-1}\otimes R^{-2})^2$, with
$\deg M=0$ and $\deg R=-2-q_0$.

The normalized conductor cover is obtained by adjoining
$\sqrt{-\Delta}$, by the quadratic equation for the ratio of
the first normal coordinate to the square of the radical
coordinate. Thus its branch divisor consists exactly of the
odd-order zeros of $\Delta$, as proved in (A.T.7). Since
$C\simeq\PP^1$ and the map has degree two, Riemann--Hurwitz
gives two distinct branch points on $H$. At each of them
$a$ vanishes: a nonzero Hessian would give the constant
transverse $A_3$ product by A.5.17 and hence an unramified
conductor cover. The coordinate comparison
(A.T.10)--(A.T.12) excludes a simple zero of $\Delta$ at
such a point. Its odd zero order is consequently at least
three. Equations (C.RA.19) now give
\[
                        A\ge2,\qquad 8-A\ge6.
\tag{C.RA.20}
\]
Therefore $A=2$. The two branch points are exactly the
zeros of $a$, each simple, and the zeros of $\Delta$ have
order three there, with no further zeros.

Every original fibre $F_j$ of (C.RA.1) gives a positive
vertical boundary component meeting $C$. The positivity
is preserved under blowup at its intersection with the
coefficient-two horizontal section, and these $r$
intersection points on $C$ are distinct. At any such point
the Hessian coefficient at its image must vanish. Otherwise
A.5.17 supplies the constant $A_3$ product, whose two
normalization branches are smooth and have no vertical
conductor; the minimal resolution $\pi$ is an isomorphism
near them. Each of the two zeros of $a$ is a branch point
of $C\to H$, and consequently has just one point above it
on $C$. It follows that $r\le2$.

Finally the exact complement calculation and (A.RT.21) give
\[
 e(S)=2-r+N_0+e(\Gamma_{\rm red})
                  \ge4-r+N_0\ge2+N_0,
\]
which proves (C.RA.18).
\end{proof}

\begin{statement}{C.2.RAT.3}{Proved}{Every uncontracted degree-one horizontal image has nonnegative Euler number}
Retain the torus-fibre, rational-normalization and ruled-boundary
morphism of C.2.RAT.1. Suppose that $C$ is not contracted and
that its actual map to the normalization of
$H=h(C)_{\rm red}$ has degree one. Let $r_{\rm dom}$ be the
number of original ruling fibres $F_j$ whose positive vertical
boundary contains a strict conductor component mapping onto
$H$. Then
\[
       0\le r_{\rm dom}\le2,\qquad
       e(S)\ge2-r_{\rm dom}\ge0.
\tag{C.RA.21}
\]
There is no assumption on the generic transverse type, on the
local complete intersection property of $H$, or on the
nilpotent structure of the conductor image.
\end{statement}
\begin{proof}
The conductor coefficient on $C$ is two, because $C$ is
the strict transform of the coefficient-two section in
(C.RA.1) and is not $\pi$-exceptional. At a general point
of $H$, its degree-one map supplies one branch of the reduced
transverse plane curve. If that branch has intrinsic
conductor exponent $c_i^{\rm int}$, its exponent in the
whole plane curve is
\[
             2=c_i^{\rm int}+\sum_{j\ne i}I(C_i,C_j).
\tag{C.RA.22}
\]
Here $I(C_i,C_j)$ is the positive intersection multiplicity
of two different plane branches. To verify the displayed
comparison, write the reduced equation as $u\prod f_j$.
On the normalization of the $i$th branch the residue
description of its dualizing generator gains the factor
$\prod_{j\ne i}f_j|_{C_i}$; its valuation is the sum in
(C.RA.22). The remaining factor is exactly the conductor
of the individual branch. This gives the actual inclusion
of its conductor ideal into its normalization, with the
stated exponent. Thus there are at most two other generic
branches. Each vertical strict component dominating $H$
supplies at least one such branch, counted with its covering
degree. This proves $r_{\rm dom}\le2$.

For an original fibre $F_j$, let $B_j$ be the positive
vertical boundary in its total transform. Its reduced
support is connected and meets $C$ at a single point.
These assertions follow at each point blowup: a centre
at the intersection of positive components inserts a
positive exceptional component; a centre at a smooth
coefficient-one point creates only a coefficient-zero
leaf; and a centre at a smooth point of coefficient
at least two creates a positive leaf. These are all the
possibilities in the exact divisor formula
$D_{\rm new}=\sigma^*D_{\rm old}-E$.
The $r$ points of $C$ so obtained are distinct.

If every component of $B_j$ is $\pi$-exceptional, there
has been at least one multiplicity-one blowup centre
above $F_j$. Indeed, in the absence of such a centre,
every component of the entire total transform of $F_j$
has positive conductor coefficient. That complete
ruling fibre, with its positive multiplicities, has
self-intersection zero. It cannot be supported on the
negative-definite exceptional matrix of a normal surface
resolution. Distinct original fibres require distinct
blowup events. Hence at most $N_0$ of the $B_j$ can be
entirely $\pi$-exceptional.

Discard those fibres and the $r_{\rm dom}$ fibres counted
above. There remain at least
\[
                        k\ge r-N_0-r_{\rm dom}
\tag{C.RA.23}
\]
fibres with a surviving strict conductor component and
none dominating $H$. Follow the connected positive
boundary from its point on $C$ to the first component
which is not $\pi$-exceptional. Its image is a curve
different from $H$ meeting $H$ at that point's image.
Every intervening connected chain is contracted by
$\pi$, so this is an actual intersection in $S$, not
an identification imposed on an auxiliary graph.

Use the normalization-incidence graph of
$G=\Gamma_{\rm red}$: its vertices are the rational
normalization components and the singular points with
more than one normalization branch, and its edges are
the individual branch preimages. The preceding $k$
intersections give $k$ distinct edges adjacent to the
vertex belonging to $H$. They remain distinct even
when two of their images are the same singular point:
$C\to H^\nu$ is an isomorphism, and the points on $C$
were distinct. Remove that vertex and all its adjacent
edges. Each of the $c$ remaining connected components
which receives one of these edges contains a
normalization vertex of a curve other than $H$.
In particular $c\le n_G-1$, where $n_G$ is the number
of irreducible components of $G$.

Contract a spanning tree in each of those $c$ graph
components. The $k$ retained edges then give a connected
multistar with $c+1$ vertices and $k$ edges. Its first
Betti number is $k-c$, so the first Betti number
$\beta$ of the full incidence graph satisfies
\[
             k-c\le\beta\le1,
          \qquad n_G\ge k+1-\beta\ge k.
\tag{C.RA.24}
\]
The upper bound on $\beta$ is the actual local torus-group
quotient (A.RT.18)--(A.RT.20). Adding genus-zero
normalization components changes no graph first Betti
number, and the normalization quotient gives
$e(G)=n_G+1-\beta\ge n_G$. Combining (C.RA.23)--(C.RA.24)
with the exact surface Euler formula yields
\[
\begin{split}
 e(S)&=2-r+N_0+e(G)\\
     &\ge2-r+N_0+n_G
      \ge2-r+N_0+(r-N_0-r_{\rm dom})
       =2-r_{\rm dom}.
\end{split}
\tag{C.RA.25}
\]
This proves the assertion.
\end{proof}

The degree-one result accounts for every generic transverse
type allowed by conductor exponent two on that branch. It
does not replace the other branches by a presumed ordinary
triple: those branches may have larger conductor exponents.
The cases in which the horizontal section is contracted, or
maps with higher degree to a singular conductor image, require
the further arguments identified separately below.

\begin{statement}{C.2.RAT.4}{Proved}{Every rational degree-two tacnode cover has positive Euler number}
Retain the torus-fibre, rational-normalization and specified
ruled-boundary morphism of C.2.RAT.1. Suppose $C$ is not
contracted, $H=h(C)_{\rm red}$ has generic transverse type
$A_3$, and the actual map $C\to H^\nu$ has degree two.
Then $e(S)\ge1$. No smoothness or local complete intersection
assumption on the image curve is required.
\end{statement}
\begin{proof}
Write $\Sigma=H^\nu\simeq\PP^1$, $G=\Gamma_{\rm red}$,
$\delta=p_a(H)\le2$, and
$s=\#\nu^{-1}(\operatorname{Sing}H)\le3$ as in
(C.RA.5)--(C.RA.6). There are exactly two generic
normalization branches above $H$, both supplied by the
degree-two map from $C$. Thus no vertical strict conductor
component dominates $H$.

We first justify using the lci jet construction without
assuming that property. The low-delta classification in the
proof of C.2.RAT.1 applies to this same reduced image curve.
Its three-independent-axis germ is excluded by
$b_1(H)\le b_1(G)\le1$. The remaining non-lci possibilities
are the space cusp with semigroup $\langle3,4,5\rangle$
and a cusp with a transverse smooth branch. The
genus-zero double-cover calculation at the end of A.5.18
excludes both: the space-cusp inequalities force
$q_0=-1$ and two simple zeros of the cleared quartic,
one over $H_{\rm reg}$, contradicting (A.T.12); for the
cusp-plus-line ring its cleared quartic has nonpositive
degree and cannot have two branch points. Therefore $H$
is lci at all its points.

Let $\Xi$ be the normalization conductor divisor of $H$,
and use the actual jet lines and sections in
(A.T.5)--(A.T.7). Put
\[
 A=\deg\operatorname{div}(a)=-2\deg Q,\qquad
 \sigma\in H^0(\Sigma,L^2(4\Xi)),\qquad
 \deg\operatorname{div}(\sigma)=8-A.
\tag{C.RA.26}
\]
Here $A$ is a nonnegative even integer and $\sigma$ is
the conductor-cleared version of $\Delta=ae-b^2/4$.
The degree-two cover $C\to\Sigma$ has two branch points,
exactly the odd-order zeros of $\sigma$. Let $B_s$ be
the number of these branch points lying over singular
points of $H$; there are $2-B_s$ over its smooth locus.
Every smooth branch point has zero order at least three
in $\sigma$, by the actual coordinate exclusion
(A.T.10)--(A.T.12).

The section $a$ must vanish at every point of $\Sigma$
over a singular point of $H$. To see this, the Hessian
matrix has entries in $\OO_H$. If its normalized scalar
$a$ is nonzero on one branch above $p$, its primitive
rank-one factor has a unit component there. A diagonal
entry consequently has a nonzero value at $p$. That
same entry has the same nonzero value on every branch,
so $a$ is nonzero at all normalization points above $p$.
In that event (A.T.13) removes the conductor denominator:
$\Delta$ is holomorphic in $L^2$ there. The orders of
$\sigma$ above $p$ are at least those of $4\Xi$,
which are even. Its total zero divisor must additionally
contain the two parity contributions of the two branch
points, whether or not they lie over $p$. Thus
\[
 \deg\operatorname{div}(\sigma)
       \ge4\deg(\Xi|_{\nu^{-1}(p)})+2
       =8\delta_p+2\ge10,
\tag{C.RA.27}
\]
contrary to (C.RA.26). This proves the claimed vanishing.

It follows in particular that
\[
                              A\ge2s-2.
\tag{C.RA.28}
\]
For $s\le1$ this is automatic; for $s=2$ the two
zeros give $A\ge2$. If $s=3$, the bounds
$\delta\le2$ and $b_1(H)\le1$ force exactly two
singular points, each of delta invariant one: one node
and one cusp. The two node preimages contribute at
least one each. At the cusp the ring is
$\mathbb C\{t^2,t^3\}$, so any vanishing Hessian entry
has order at least two. Dividing such a diagonal entry
by the square of a unit component of the primitive
normalized factor shows that $a$ has order at least
two there. Thus $A\ge1+1+2=4=2s-2$.

We now check exactly which smooth image points can
receive two vertical boundary attachments on their two
conductor preimages. At a smooth point with $a\ne0$,
A.5.17 gives the constant $A_3$ product and no vertical
attachment. Suppose instead that $a=0$ and $\Delta$
is a unit. In the actual coordinates (A.T.10), the
coefficient $b$ is then a unit. Substitute $x=y^2\lambda$
in that equation and divide by $y^4$. The result is
holomorphic in $(\lambda,y,z)$ and has, at the point
in question, the simple root
$\lambda=-e(0)/b(0)$. The implicit-function theorem
gives a holomorphic $\lambda(y,z)$ and an actual
factorization
\[
 f_0=F_1F_2,\qquad
 F_1=x-y^2\lambda(y,z),\qquad
 F_2|_{F_1=0}=y^2\,u(y,z),\quad u(0)\ne0.
\tag{C.RA.29}
\]
The last equality is the restriction of
$\partial f_0/\partial x$ to $F_1=0$; its quotient
by $y^2$ has constant value $b(0)$.
The branch $F_1=0$ is smooth, and its conductor in the
whole reduced hypersurface is exactly the divisor
$2\{y=0\}$, by the intersection factor in
(C.RA.22). There is therefore no vertical conductor
near this normalization point. Its minimal resolution
is an isomorphism there as well. Since $\Delta$ is a
unit, this is one of the two distinct conductor
preimages over the given point of $\Sigma$.
Consequently a smooth nonbranch value with two vertical
attachments has a positive even zero order of $\sigma$,
at least two.

Let $d_{\rm reg}$ count distinct smooth values with
two such attachments. The two branch points and these
distinct nonbranch zeros give
\[
             3(2-B_s)+B_s+2d_{\rm reg}\le8-A.
\tag{C.RA.30}
\]
Among the $s$ points above the singular locus, the
$B_s$ branch points each have only one preimage on
$C$. Thus at most $s-B_s$ of these values can have
two attachments. If $d$ counts all values on
$\Sigma$ with two attachments, (C.RA.28)--(C.RA.30)
give
\[
 d\le s-B_s+d_{\rm reg}
       \le s+1-\frac A2\le2.
\tag{C.RA.31}
\]

In fact $d\le1$. If $d=2$, (C.RA.31) requires
$A\le2s-2$. For $s=0$ this is impossible, and for
$s=1$ it contradicts the vanishing of $a$ over the
singular point. For $s=2$, it forces $A=2$ and
the two singular preimages exhaust the zero divisor
of $a$. For $s=3$, it forces $A=4$, exhausted by
the node and cusp orders $1+1+2$ established above.
In both remaining cases $a$ is nowhere zero over
$H_{\rm reg}$. The constant $A_3$ product there
has neither a cover branch point nor a vertical
attachment. Thus both cover branch points lie
over $\operatorname{Sing}H$, and only the other
$s-2\le1$ singular preimages can have two
attachments. This again contradicts $d=2$.
Consequently
\[
                            d\le1.
\tag{C.RA.31a}
\]

It remains to connect this local count to the exact
surface Euler number. As in C.2.RAT.3, at least
$K\ge r-N_0$ original positive vertical boundaries
have a component which is not $\pi$-exceptional.
None maps onto $H$, so they give $K$ distinct points
of $C$ anchored to image curves other than $H$.
After projection to $\Sigma$ their number of distinct
normalization-incidence edges is $k=K-d'$ where
$d'\le d\le1$: a degree-two fibre has at most two
points, and the loss of one occurs precisely at a
value with two selected attachments. Apply the graph
construction of (C.RA.24) to these $k$ distinct
edges. It gives $n_G\ge k$ and $e(G)\ge n_G$.
The actual Euler decomposition therefore yields
\[
\begin{split}
 e(S)&=2-r+N_0+e(G)
       \ge2-K+e(G)\\
     &\ge2-K+k=2-d'\ge1.
\end{split}
\tag{C.RA.32}
\]
Every contraction, covering fibre and graph edge in
this count comes from the specified maps
$Z\to T\to S$ and $C\to\Sigma\to H$.
\end{proof}
